\documentclass[10pt,a4paper]{amsart}
\usepackage[margin=19mm]{geometry}
\usepackage{amsmath,amssymb,mathtools}
\usepackage[scr=rsfs,cal=euler]{mathalfa}
\usepackage{xfrac}
\usepackage[unicode,hidelinks,bookmarks=false]{hyperref}
\newcommand{\ie}{i.e.\ }
\newcommand{\cf}{cf.\ }
\newcommand{\resp}{respectively\ }
\newcommand{\Subsection}{\subsection}
\providecommand{\nobreakdash}{\nobreak\hbox{-}}
\setlength{\emergencystretch}{2em}
\multlinegap0pt
\usepackage{amssymb}
\usepackage{xcolor}
\usepackage{tikz}
\newtheorem{theorem}{Theorem}
\newcommand{\size}[1]{\left|#1\right|}
\title{\textbf{A note on vertex-critical\\induced subgraphs of shift graphs}}
\author{Tom\'a\v{s} Kaiser, Mat\v{e}j Stehl\'ik, Riste \v{S}krekovski, Bjarne Toft}
\date{Source: arXiv:2601.16342v2 (CC BY 4.0)}
\begin{document}
\maketitle

\noindent\emph{Source and licence.} \textbf{A note on vertex-critical\\induced subgraphs of shift graphs}. Version arXiv:2601.16342v2 (CC BY 4.0), distributed by arXiv under the Creative Commons Attribution 4.0 International licence, http://creativecommons.org/licenses/by/4.0/, as recorded in the arXiv licence metadata for this exact version and shown on the abstract page https://arxiv.org/abs/2601.16342v2. Full licence notice: \texttt{licenses/arxiv-2601.16342-CC-BY-4.0.txt}.\par\smallskip

\noindent\textit{Editorial scope.} Counted unit: the article's theorem t:crit-all (source0.tex lines 204--206). The article's complete original proof of it is delivered with it (source0.tex lines 207--262, one \texttt{proof} environment of 56 lines). No mathematics is written, repaired or re-derived: the statement and the proof are the article's own lines, in the article's own notation, and no retained line is substituted or re-flowed.  No further source-local context is needed: every cross-reference of every retained line resolves inside the two delivered spans.  Nothing was omitted from any retained span, so the package carries no omission record.  No retained line contains a citation, so the wrapper carries no bibliography and no bibliography entry.  No retained line contains a figure environment, an image-inclusion command or drawing code, so no image is delivered and the package declares no asset dependency.  Editorial formatting only, in addition: an AMS \texttt{amsart} preamble that restates the article's own declarations and macros used by the retained lines, namely the environments \texttt{theorem} on separate counters, together with the standard packages those lines need.  The retained environments print in this document as Theorem 1 in order of appearance, whereas the article prints the same items with the numbering of its own full document.  The complete native source of the article as distributed in the arXiv e-print for this exact version is bundled unchanged as \texttt{sources/193230/source0.tex}.
\begin{theorem}\label{t:crit-all}
  $\chi(G[W]) = \chi(G) = n+1.$
\end{theorem}

\begin{proof}
  By contradiction. Consider a $W$-good sequence
  $a = (a_1,\dots,a_{2^n+1})$. We first transform the sequence to another
  $W$-good sequence using the following procedure. If there is an
  index $r$ such that some proper subset $a'_r$ of $a_r$ is not
  included among $a_{r+1},\dots,a_{2^n+1}$, let $s$ be the largest
  such index and replace $a_s$ (at position $s$) in the sequence with
  $a'_s$. We prove that the resulting sequence is $W$-good. Suppose
  not. Since $a'_s\subseteq a_s$, there cannot be $i < s$ such that
  $(i,s)\in W$ and $a_i\subseteq a'_s$ as this would contradict the
  $W$-goodness of the original sequence. Thus, let us consider $j > s$
  such that $(s,j)\in W$. In this case, $a'_s$ cannot be a subset of
  $a_j$, because it would be a proper subset (by the choice of $s$),
  so by the maximality of $s$, it would be included among
  $a_{j+1},\dots,a_{2^n+1}$, contradicting the choice of $s$. Thus,
  the resulting sequence is $W$-good, as claimed.

  Relabel the new element $a'_s$ as $a_s$, and repeat the above step
  until this is no longer possible. Clearly, the
  algorithm only takes a finite number of steps and terminates with a
  $W$-good sequence $a_1,\dots,a_{2^n+1}$ with the following property:
  \begin{itemize}
  \item[] For each $i\in [1,2^n+1]$ and every proper subset $b \subset a_i$, there is $j \in [i+1,2^n+1]$ such that $a_j = b$.
  \end{itemize}
 
  Observe that, as a
  consequence, $a_{2^n+1} = \emptyset$. Furthermore, for each
  $j\in[1,2^n+1]$,
  \begin{equation}
    \label{eq:1}
    \text{if }\size{a_j} = k, \text{ then } j \leq 2^n - 2^k + 2,
  \end{equation}
  as there has to be space to the right of $a_j$ for $2^k-1$ proper
  subsets of $a_j$.

  Since the number of subsets of $[1,n]$ is $2^n$ while the length of the sequence
  $a$ is $2^n+1$, one of the subsets needs to occur multiple times in $a$. In particular, there exists a pair $i,j\in [1,2^n+1]$
  such that:
  \begin{itemize} 
    \item[] $i < j$ and $a_i \subseteq a_j$; moreover, we choose this pair
  such that $i$ is as large as possible and, subject to this condition, $j$ is as small as possible.
  \end{itemize}
  Let us write $k = \size{a_j}$. 

  By the maximality of $i$, the subsequence $(a_{i+1}, \dots, a_{2^n+1})$ of $a$ consists of pairwise distinct subsets of $[1,n]$. In addition, it is easy to see that by the choice of $i$ and $j$, the only superset of $a_j$ in this subsequence is $a_j$ itself. The fact that the remaining $2^{n-k}-1$ supersets are missing implies an upper bound on the length of the subsequence:
  \begin{equation*}
    2^n - i + 1 \leq 2^n - (2^{n-k}-1),
  \end{equation*}
  from which we find $2^{n-k} \leq i$. In combination with~\eqref{eq:1}, we obtain
  \begin{equation}
    \label{eq:6}
    2^{n-k} \leq i < j \leq 2^n - 2^k + 2.
  \end{equation}
  This just says that $i,j\in I_{n-k}$. As a consequence, $(i,j)\in
  W$, which contradicts the $W$-goodness of $a$ and concludes the proof of Theorem~\ref{t:crit-all}.
\end{proof}

\end{document}
